\section{\titleUtransformRom}\label{sec:utransform}
  
  We will now meet the Fujisaki--Okamoto (FO) transformation~\cite{C:FujOka99},
  	a PKE-to-KEM recipe that became quite popular in the post-quantum-context --
  	it is used as part of many post-quantum KEMs,
  	including the now-standardized ML-KEM (Kyber)~\cite{FIPS203}. 
  
  We first focus on the part of FO that turns a deterministic PKE scheme into a KEM.
	  Following the literature~\cite{TCC:HofHovKil17}, this part of FO will be called the \U-transform.
	  %, as defined by Hofheinz et al. (hereafter HHK).
	  To have a simple ROM example, we will -- for now -- only prove that the obtained KEM
	  achieves key indistinguishability against \emph{passive} adversaries.
	  This proof will already teach us several useful tricks you can play in the ROM.
	  These tricks foreshadow and motivate the QROM tools that will be introduced in~\cref{sec:o2h}
	  \todo{\cref{sec:o2h} vs \cref{sec:C-O2H}}.

  At a later point\todo{add pointer to where},
    we will also prove key indistinguishability against \emph{active} adversaries.
    There also exists a second part of FO that turns probabilistic PKE schemes
    into ones that fit the \U-transform's requirements.
    (This part is called the \T-transform, in case you're curious, see App.~\ref{sec:ttransform}.)
%    (\T) de-randomizes it by first deriving the random coins by hashing
%    the message (and possibly the public key), then running the encryption
%    algorithm using those coins; then, after decrypting, it re-derives the
%    random coins and re-encrypts the message to verify that resulting
%    ciphertext matches the input. 

  In short, the \U-transformation builds a KEM from a PKE scheme by sampling a random message,
    encrypting it, and hashing it (possibly together with the public key and ciphertext) to derive a
    symmetric key $\eK$. This is the transformation that will serve as our
    running example.\footnote{
      When Hofheinz et al.~\cite{TCC:HofHovKil17} first defined
      the \T- and \U-transforms, they actually defined several \U-variants;
      %that differ from the variant discussed here
  		%in how they handle invalid ciphertexts and in what they hash into the
      %output key; 
      \cref{def:U-transform} corresponds to
      \UExpMess in their notation (message-only hashes, explicit decryption failures).
    }

  \begin{definition}[\U-transform]\label{def:U-transform}
    Let $\dPKE = (\KG, \Enc, \Dec)$ be a deterministic public-key encryption scheme, and let 
    $\ROkey: \bin^* \to \bin^m$ be a hash function.
    Then the \U-transform is the key-encapsulation mechanism 
    $\U[\PKE, \ROkey] \eqqcolon (\KG, \UEncaps, \UDecaps)$, as defined in \cref{fig:def:U-transform}.
    \begin{figure}[ht]
   	  \begin{center}
  		  \fbox{
  		  	\small
  		  	\nicoresetlinenr
  		  	%\begin{minipage}[t]{3.3cm}
  		  	%	\underline{$\UKG()$}
  		  	%	\begin{nicodemus}
  		  	%		\item $(\pk, \sk) \sample \KG()$
  		  	%		%\item $\rejectionSeed \sample \MSpace$
  		  	%		%\item $\sk' \leftarrow (\sk, \rejectionSeed)$
  		  	%		\item \pcreturn $(\pk, \sk)$
  		  	%	\end{nicodemus}
  		  	%\end{minipage}
  		  	%
  		  	%\quad
  		  	
  		  	\begin{minipage}[t]{3cm}
            \underline{$\UEncaps_\pk()$}
  		  		\begin{nicodemus}
              \smallskip
  		  			\item $x \sample \bin^n$
              \smallskip
  		  			\item $c \gets \Enc(\pk, x)$
              \smallskip
  		  			\item $\eK \gets \ROkey(x)$
              \smallskip
  		  			\item \pcreturn $(\eK, c)$
  		  		\end{nicodemus}
  		  	\end{minipage}
  		  	
  		  	\quad
  		  	
  		  	\begin{minipage}[t]{2.5cm}
           \underline{$\UDecaps_\sk(c)$}
  		  		\begin{nicodemus}
              \smallskip
  		  			%\item Parse $\sk' := (\sk, \rejectionSeed)$
              \item $x' \gets \Dec_\sk(c)$
              \smallskip
  		  			\item \pcif $x' = \bot :$ 
              \smallskip
              \item   \quad \pcreturn $\bot$
              \smallskip
  		  				%\item \pcreturn $\ROkey(\rejectionSeed, c)$
  		  			\item \pcreturn $\ROkey(x')$
  		  		\end{nicodemus}
  		  	\end{minipage}
  		  }
  	  \end{center}
  	  \caption{
  	    \label{fig:def:U-transform}			
         Algorithms of $\U[\PKE, \ROkey]$.
      }
    \end{figure} 
  \end{definition}

  %\hhnote{
  %  Leave this version without the extra re-encryption check (which also
  %  matches the original \U), and define a new \U with the re-encryption check
  %  when we later want CCA security.
  %}

  As explained in \cref{sec:basics}, One-way security under Chosen
    Plaintext Attacks, or \owcpa, is really a minimal security
    requirement for PKE schemes, and should be considered too weak a security
    goal for most use-cases. 
    %The goal of One-Wayness (OW) simply says that
    %given a ciphertext, it should be hard to recover the plaintext, and the
    %power of Chosen Plaintext Attacks (CPA), the ability to encrypt any
    %message under the key you are targetting, are of course ever-present when
    %encryption is public.
    Much preferred would be the stronger security requirement of
    \emph{indistinguishability} under chosen-plaintext attacks, or
    \indcpa.\footnote{
      The standard security requirement for deployable KEMs is an even stronger one,
      namely indistinguishability under chosen-\emph{ciphertext} attacks, or
      \indcca. 
      As it turns out, the \U-transform also achieves this, provided the
      re-encryption check from the \T-transform is carried forward; this will
      be covered in later chapters (to appear).
    } 

  However, stronger security assumptions usually mean less efficient
    schemes, as the parameters must be increased to compensate for the
    increased attack surface. What's worse, we are taking as our starting
    point a \emph{deterministic} encryption scheme. Such a PKE scheme could
    \emph{never} achieve indistinguishability, since the adversary can simply
    re-encrypt the two candidate messages, and check which of them
    encrypts to the challenge ciphertext.

  And yet, taking our deterministic PKE scheme through the \U-transform
    does in fact give us a KEM that achieves indistinguishability
    \ldots and the object that allows us to cross that chasm
    %between one-wayness and indistinguishability 
    is precisely the \emph{random oracle}.

  The proof strategy will rely on the fact that our reduction can
    fully control the random oracle for the adversary, inspecting the queries
    before answering them via the lazy-sampling procedure from \cref{sec:lazy-sampling}. 
    Intuitively speaking, in order to successfully distinguish, the adversary
    will have to have called the random oracle on one of
    the messages used to derive the ephemeral key. 
    The reduction can recognize this the moment it happens, by re-encrypting
    each message and checking if the encryption matches the challenge it was
    issued; if so, it wins at the one-wayness game by outputting that message.

    %able to recognize the correct solution; this fact allows us to get
    %a tight bound in this case. (Otherwise we would have to guess which of the
    %calls to the random oracle contained the solution we are looking for,
    %which would incur a tightness loss related to the probability that
    %we guess correctly.)

  To summarize, if $\dPKE$ satisfies one-wayness
    (\cref{def:owcpa}), then $\KEM$ satisfies indistinguishability
    (\cref{def:indcca}), as captured in the following theorem.

  \begin{theorem}[\dPKE is \owcpa $\To$ \KEM is \indcpa]
    \label{thm:cparom}
    Let \dPKE be a deterministic PKE scheme, let \RO be a random oracle, 
    and let the $\KEM = \U[\dPKE, \RO]$.
    Then there is a simple, fully black-box reduction $\bdv$, such that for
    every adversary $\adv$,
    %that makes at most $\qe$ calls to the challenge oracle, 
    \[
      \advantage{\indcpa}{\KEM}[\left(\adv^{\RO}\right)] 
        \leq \advantage{\owcpa}{\dPKE}[\left(\bdv^{\adv}\right)]\, .
    \]
    The adversary $\bdv$ runs in the same time as $\adv$, plus the %(relatively small) 
    overhead of lazy-sampling a random oracle.
  \end{theorem}

  \begin{proof} %[\cref{thm:cparom}]
    We want to construct a $\bdv$ that exploits any winning
      strategy of $\adv$ in such a way that if $\adv$ \emph{would} win at the
      indistinguishability game, then $\bdv$ \emph{actually} wins
      at the one-wayness game.
      However, the constructed adversary $\bdv$ will not be able simulate the
      real ($b=0$) world of the KEM indistinguishability game, since it does
      not know which message the challenge ciphertext encrypts -- that's what
      it is trying to figure out! 
      %hence, it can not hash the message to produce the
      %correct symmetric key for $\adv$.
      It can, however, easily simulate the random ($b=1$) world, since here the
      symmetric key given to the adversary is simply uniformly sampled.

    But in order to argue that $\adv$ can not distinguish $\bdv$'s simulation
      from the real-world game, we need to alter the game in such a way that the
      move from real to random keys is completely hidden from the view of $\adv$.
      We do this by inserting two intermediary games between the real-world and the
      random-world games, as illustrated in \cref{fig:cparom-hops-illustrated}.
      These two middle games formalize the intuition that, thanks to the \U-transform,
      the adversary will have to call the random oracle on the correct message
      in order to notice that it is not in the real world: If it
      does, then a boolean flag $\bad$ is set to true. Assuming the $\bad$ flag is
      never set to true, the game proceeds exactly as the real-world game, and
      thanks to \cref{lemma:bad}, we know that the noticeable difference
      between these games is bounded exactly by the probability that $\bad$ is
      set to true. Then, \emph{conditioned} on $\bad$ remaining false,
      we can exchange the real symmetric key for a randomly
      sampled one, and be sure that the adversary will never be able to
      notice.
      %Therefore, the second game hop from real to random keys is ``for free'',
      %since the two games are equal in terms of probabilities. 


      \crakzfigcode{cparom-hops-illustrated}{
        Proof outline for \cref{thm:cparom}. The arrows show the changes
        incurred by each game hop, and the braces show the difference in
        win probabilities between each game.
        %Here, $\indcpa_0$ resp.\ $\indcpa_1$ is shorthand for
        %the game $\experiment{\indcpa}{\pkescheme}[]$ (\cref{fig:ind-ow}) instantiated with challenge bit
        %$\chalbit = 0$ (left world) resp.\ $\chalbit = 1$ (right world).
      }
  
    The four games are given in pseudocode in \cref{fig:cparom-hops}, in a
      layout matching \cref{fig:cparom-hops-illustrated}, and with
      changes between games highlighted.
      To make the argument formal, start by recalling the definition of a
      distinguishing advantage; we will use the ``difference'' version of
      distinguishing advantages, as
      it fits more cleanly with this ``game hopping'' style of proof.
      (For the sake of simplicity, we suppress the superscript dependencies
      $\RO$ and $\adv$ in the following.)
      \[
      	\advantage{\indcpa}{\KEM} \coloneqq \abs{
      		\condprobsample{
      			\experiment{\indcpa}{\KEM}
      		}{1 \gets \adv}{\chalbit = 0}
      		- \condprobsample{
      			\experiment{\indcpa}{\KEM}
      		}{1 \gets \adv}{\chalbit = 1}
      	}\, .
      \]
      The game $\game 1$ is just the real-world ($\chalbit=0$) $\indcpa$ game, with the
      construction of \KEM written out explicitly, together with a lazy-sampled
      random oracle. Similarly, $\game 4$ is just the random-world
      ($\chalbit=1$) $\indcpa$ game with the \KEM construction explicit. We
      can therefore write
      \[
      	\advantage{\indcpa}{\KEM} = \abs{
      		\probsample{ \game{1}(\adv) }{1 \gets \adv}
      		- \probsample{ \game{4}(\adv) }{1 \gets \adv}
      	}\, .
      \]
      We are of course always free to ``add zero'' by adding and subtracting
      the same term, so this is equal to:
      \begin{align*}
        \advantage{\indcpa}{\KEM} &= 
          |\probsample{ \game{1}(\adv) }{1 \gets \adv} 
            - \probsample{ \game{2}(\adv) }{1 \gets \adv} \\
          &\qquad + \probsample{ \game{2}(\adv) }{1 \gets \adv} 
            - \probsample{ \game{3}(\adv) }{1 \gets \adv} \\
          &\qquad + \probsample{ \game{3}(\adv) }{1 \gets \adv} 
            - \probsample{ \game{4}(\adv) }{1 \gets \adv}|\, .
      \end{align*}
      Applying the triangle inequality then gives us an upper-bound 
        on the original distinguishing advantage as the sum of the distinguishing
        advantages between successive games:
      \begin{align*}
        \advantage{\indcpa}{\KEM} &\leq 
          \abs{\probsample{ \game{1}(\adv) }{1 \gets \adv} 
            - \probsample{ \game{2}(\adv) }{1 \gets \adv}} \\
          &\qquad + \abs{\probsample{ \game{2}(\adv) }{1 \gets \adv} 
            - \probsample{ \game{3}(\adv) }{1 \gets \adv}} \\
          &\qquad + \abs{\probsample{ \game{3}(\adv) }{1 \gets \adv}
            - \probsample{ \game{4}(\adv) }{1 \gets \adv}}\, .
      \end{align*}

      Let us analyze each distinguishing advantage in turn.
        First, instead of deriving the key from the hash of the message as
        usual, $\game 2$ instead samples $\eK^*$ uniformly at random at the start
        of the game and gives it to $\adv$. Then, if $\RO$ is ever called with
        the encrypted message $x^*$, which \emph{should} have been hashed to
        derive $\eK^*$ at the start of the game, the oracle database is updated such that
        $\Db[x^*] = \eK^*$. Since the random oracle was lazy-sampled anyway,
        there is no way for $\adv$ to tell if the internal oracle database was
        updated right away, or only when it is needed. From the point of view
        of $\adv$, these games therefore behave identically, and
      \begin{align*}
        \abs{\probsample{ \game{1}(\adv) }{1 \gets \adv} 
          - \probsample{ \game{2}(\adv) }{1 \gets \adv}} = 0\, . 
      \end{align*}
      Two more things to notice about $\game 2$: First, the oracle interface
        could have simply done the above check by checking whether the input $x$
        matches the sampled message $x^*$; instead, it takes the roundabout
        route of re-encrypting the input and checking whether the resulting
        ciphertext is equal to $\ctxt^*$. Since \dPKE is deterministic
        (and by assumption perfectly correct), these checks lead to identical
        outcomes; but of course, only the latter check will be available to $\bdv$,
        who does not know $x^*$. Second, in addition to the ``delayed
        programming'' of the random oracle, the game introduces the $\bad$
        flag in line $10$. Initialized to $\false$, this flag is only set to
        $\true$ in the case that the database is actually updated to hold the
        sampled $\eK^*$, or in other words, if the random oracle is called on
        the encrypted message $x^*$. Other than this, the flag has no effect on the
        behaviour of the game.

      Compared to the changes introduced in $\game 2$, the change from 
        $\game 2$ to $\game 3$ is deceptively simple: When the random oracle is
        called on $x^*$, after $\bad$ is set to $\true$, do \ldots nothing.
        This is where $\game 2$ would program the random oracle with $\eK^*$,
        thus ensuring that the random oracle remained consistent with the
        real world. Now we are removing this programming, which means that
        once $\bad$ is set to true,
        the game is \emph{noticably} inconsistent with the real-world game:
        If $\adv$ was able to recover $x^*$ from $\ctxt^*$, then they will
        quickly notice that $\RO(x^*)$ does not evaluate to $\eK^*$ (except
        if, by freak accident, the random value drawn by the lazy-sampling
        procedure happens to be exactly $\eK^*$).

      Observe how $\game 2$ and $\game 3$ are by design exactly identicaly
        up until the point that $\bad$ is set to $\true$. This means that we
        can use \cref{lemma:bad} to bound their difference as the probability
        that $\bad$ is set to $\true$ at the end of either game:
      \begin{align*}
        \abs{\probsample{ \game{1}(\adv) }{1 \gets \adv} 
          - \probsample{ \game{2}(\adv) }{1 \gets \adv}} \leq
          \probsample{\game{2}(\adv)}{\bad} = \probsample{\game{3}(\adv)}{\bad}\, . 
      \end{align*}

    %%% Game hops, single-challenge version
    \begin{figure}[t]
    	\begin{center}
    	\resizebox{\textwidth}{!}{%
    	\begin{tabular}{@{}c@{\hspace{0em}}c@{}}
        % TOP LEFT: Game 1
    		\fbox{
    			\small
    			\nicoresetlinenr
          \begin{minipage}[t][3.4cm]{5cm}	
            \underline{\gameHeading $\game 1(\adv)$}
            \smallskip
    				\begin{nicodemus}
              \item $(\pk, \sk) \gets \KG()$
              \smallskip
              \item $x^* \sample \bin^n$
              \smallskip
              \item $\ctxt^* \assign \Enc_{\pk}(x^*)$
              \smallskip
              \item $\eK^* \gets \RO(x^*)$ 
              \smallskip
              \item $\advbit \gets \adv^{\RO}(\pk, \eK^*, \ctxt^*)$ 
              \smallskip
              \item $\pcreturn \advbit$
    				\end{nicodemus}
    			\end{minipage}
    
          \begin{minipage}[t]{3.6cm}
    				\underline{\oracleHeading $\RO(x)$}
            \smallskip
    				\begin{nicodemus}
              \item $\pcif \Db[x] = \bot :$ 
              \smallskip
              \item   \quad $\Db[x] \sample \bin^m$
              \smallskip
              \item $\pcreturn \Db[x]$
    				\end{nicodemus}	
    			\end{minipage}
    		}
        &
        %% TOP RIGHT: Game 4
    		\fbox{
    			\small
    			\nicoresetlinenr
          \begin{minipage}[t][3.4cm]{4.2cm}	
            \underline{\gameHeading $\game 4(\adv)$}
            \smallskip
    				\begin{nicodemus}
              \item $(\pk, \sk) \gets \KG()$
              \smallskip
              \item $x^* \sample \bin^n$
              \smallskip
              \item $\ctxt^* \assign \Enc_{\pk}(x^*)$
              \smallskip
              \item \dbox{$\eK^* \sample \bin^m$}
              \smallskip
              \item $\advbit \gets \adv^{\RO}(\pk, \eK^*, \ctxt^*)$ 
              \smallskip
              \item $\pcoutput \advbit$
    				\end{nicodemus}
    			\end{minipage}
    
    			\begin{minipage}[t]{3.6cm}
    				\underline{\oracleHeading $\RO(x)$}
            \smallskip
    				\begin{nicodemus}
              \item \sout{$\ctxt \assign \Enc_\pk(x)$}
              \smallskip
              \item \sout{$\pcif \ctxt = \ctxt^* :$}
              \item \quad \sout{$\red{\bad \assign \true}$}
              \smallskip
              \item $\pcif \Db[x] = \bot :$ 
              \smallskip
              \item   \quad $\Db[x] \sample \bin^m$
              \smallskip
              \item $\pcreturn \Db[x]$
    				\end{nicodemus}	
    			\end{minipage}
    		}
        \\[2ex]
        % BOTTOM LEFT: Game 2
    		\fbox{
    			\small
    			\nicoresetlinenr
          \begin{minipage}[t][4.2cm]{5cm}	
            \underline{\gameHeading $\game 2(\adv)$}
            \smallskip
    				\begin{nicodemus}
              \item \fbox{$\bad \assign \false$}
              \smallskip
              \item $(\pk, \sk) \gets \KG()$
              \smallskip
              \item $x^* \sample \bin^n$
              \smallskip
              \item $\ctxt^* \assign \Enc_{\pk}(x^*)$
              \smallskip
              \item \sout{$\eK^* \gets \RO(x^*)$} \fbox{$\eK^* \sample \bin^m$}
              \smallskip
              \item $\advbit \gets \adv^{\RO}(\pk, \eK^*, \ctxt^*)$ 
              \smallskip
              \item $\pcoutput \advbit$
              \smallskip
    				\end{nicodemus}
    			\end{minipage}
    
    			\begin{minipage}[t]{3.6cm}
    				\underline{\oracleHeading $\RO(x)$}
            \smallskip
    				\begin{nicodemus}
              \item \fbox{$\ctxt \assign \Enc_\pk(x)$}
              \item \fbox{$\pcif \ctxt = \ctxt^* :$} 
              \item \quad \fbox{$\red{\bad \assign \true}$} 
              \item \quad \fbox{$\Db[x] \assign \eK^*$} 
              \smallskip
              \item $\pcif \Db[x] = \bot :$ 
              \smallskip
              \item   \quad $\Db[x] \sample \bin^m$
              \smallskip
              \item $\pcreturn \Db[x]$
    				\end{nicodemus}	
    			\end{minipage}
    		}
        &
        % BOTTOM RIGHT: Game 3
    		\fbox{
    			\small
    			\nicoresetlinenr
          \begin{minipage}[t][4.2cm]{4.2cm}	
            \underline{\gameHeading $\game 3(\adv)$}
            \smallskip
    				\begin{nicodemus}
              \item $\bad \assign \false$
              \smallskip
              \item $(\pk, \sk) \gets \KG()$
              \smallskip
              \item $x^* \sample \bin^n$
              \smallskip
              \item $\ctxt^* \assign \Enc_{\pk}(x^*)$
              \smallskip
              \item $\eK^* \sample \bin^m$
              \smallskip
              \item $\advbit \gets \adv^{\RO}(\pk, \eK^*, \ctxt^*)$ 
              \smallskip
              \item $\pcoutput \advbit$
    				\end{nicodemus}
    			\end{minipage}
    
    			\begin{minipage}[t]{3.6cm}
    				\underline{\oracleHeading $\RO(x)$}
            \smallskip
    				\begin{nicodemus}
              \item $\ctxt \assign \Enc_\pk(x)$
              \smallskip
              \item $\pcif \ctxt = \ctxt^* :$
              \item \quad $\red{\bad \assign \true}$
              \item \quad \sout{$\Db[x] \assign \eK^*$} 
              \smallskip
              \item $\pcif \Db[x] = \bot :$ 
              \smallskip
              \item   \quad $\Db[x] \sample \bin^m$
              \smallskip
              \item $\pcreturn \Db[x]$
    				\end{nicodemus}	
    			\end{minipage}
    		}
    	\end{tabular}%
    	}
    	\end{center}
    	\caption{\label{fig:cparom-hops}
        The game hops in the proof of \cref{thm:cparom}, following
        the arrows of \cref{fig:cparom-hops-illustrated}.
        The dashed box highlights the difference between the real-world
        and the random-world indistinguishability games, the solid boxes
        highlight lines of code that were added between successive games, and
        strike-outs highlight lines of codes that were removed between games.
        Highlighted in \red{red} is the ``bad event'', represented by the flag
        $\bad$ being set to $\true$.
      }
    \end{figure}

      So, $\game 3$ will be noticeably inconsistent with the real-world game if
        $\adv$ ever calls $\RO$ on $x^*$; but what $\game 3$ \emph{will} look
        consistent with, is the case that $\adv$ was given a uniformly random
        $\eK$, chosen completely independently of the ciphertext $\ctxt^*$.
        That's the random-world game! 
        The final hop to $\game 4$ makes this manifest by removing the
        the $\bad$ flag again, as well as the re-encryption check in the
        oracle interface. Since the flag itself had no affect on the operation
        of the game, this change is once again for free, and the result is
        exactly the random-world $\indcpa$ game instantiated with our
        construction.
      \begin{align*}
        \abs{\probsample{ \game{3}(\adv) }{1 \gets \adv} 
          - \probsample{ \game{4}(\adv) }{1 \gets \adv}} = 0\, . 
      \end{align*}

      Collecting, we see that we have found the bound
      \[
        \advantage{\indcpa}{\KEM} 
            %\abs{\probsample{ \game{1}(\adv) }{1 \gets \adv} 
            %- \probsample{ \game{2}(\adv) }{1 \gets \adv}} \\
            %+ \abs{\probsample{ \game{2}(\adv) }{1 \gets \adv} 
            %- \probsample{ \game{3}(\adv) }{1 \gets \adv}}
            %+ \abs{\probsample{ \game{3}(\adv) }{1 \gets \adv}
            %- \probsample{ \game{4}(\adv) }{1 \gets \adv}} \\
            \leq 0 + \probsample{\game{3}(\adv)}{\bad} + 0 =
            \probsample{\game{3}(\adv)}{\bad}\, .
      \]

      All that remains is to bound $\probsample{\game{3}(\adv)}{\bad}$,
        and this is where our reduction $\bdv$ enters the stage.
        The reduction is playing its own $\owcpa$ game, and is given black-box
        access to the adversary $\adv$ (presumably together with high praises
        of its code-breaking abilities). It wants to leverage the abilities of
        $\adv$ to win at its own game, by simulating an environment for
        $\adv$ and extracting its outputs. However, $\bdv$ will have to
        provide an \emph{exact} simulation of the environment that $\adv$
        expects: if $\adv$ notices anything is off, it can refuse to
        participate, for example by deliberately outputting wrong values, or
        halting prematurely.\footnote{
          To see how this admittedly anthropomorphizing description arises from the
          mathematical statement, recall that the theorem quantizes
          over the distinguishing advantage of \emph{any} adversary that is
          placed in the \indcpa environment. It does not (and can not) say anything
          about how these same adversaries behave in any \emph{other}
          environments, and the ``for all'' quantifier means that we 
          might as well assume the worst.
        }

      The problem is, $\adv$ expects to be placed in the real- or random-world
        $\indcpa$ game with 50/50 probability, but $\bdv$ only knows how to
        simulate the random-world game. 
        If $\bdv$ just simulated the random-world game for $\adv$,
        then it would look to $\adv$ like $b$ were equal to $1$ with $100\%$
        probability. Hence, it would seem that an adversary that simply ignores
        its input and outputs $1$ would break $\bdv$'s strategy, since it
        would win at the simulated game $100\%$ of the time while leaving $\bdv$
        empty-handed.

      But of course, $\bdv$ will not be simulating the random-world game
        $\game 4$, but rather the equivalent $\game 3$. Notice how the $\bad$
        flag is first set to $\true$ the first time the encrypted value $x^*$
        is input to the random oracle. Notice further that, thanks to the
        re-encryption trick, this check can actually be implemented by $\bdv$,
        who knows $\pk$ and $\ctxt^*$ but not $x^*$. A strategy reveals
        itself: Use the inputs $\pk$ and $\ctxt^*$ given by the \owcpa
        environment to simulate $\game 3$ for $\adv$, except that when $\game
        3$ would set $\bad$ to $\true$, instead halt the simulation and output
        the value $x^*$ that was just entered into the random oracle.

      The pseudocode for our reduction $\bdv$ is given in
        \cref{fig:cparom-red}. Since the event that $\bdv$ halts and outputs
        the correct value corresponds exactly to the event that $\game 3$
        would set $\bad$ to $\true$, we have
      \[
        \probsample{\game{3}(\adv)}{\bad} =
          \advantage{\owcpa}{\dPKE}[(\bdv)]\, ,
      \]
      from which we may finally conclude that
      \[
        \advantage{\indcpa}{\KEM} \leq 
        \probsample{\game{3}(\adv)}{\bad}
        = \advantage{\owcpa}{\dPKE}[(\bdv)]\, . 
      \]

    %%% Reduction
    \begin{figure}[t]
    	\begin{center}
    		\fbox{
    			\small
    			\nicoresetlinenr
    			\begin{minipage}[t]{3.7cm}	
            \underline{\redHeading $\bdv(\pk, \chalctxt)$}
            \smallskip
    				\begin{nicodemus}
              \item $\eK^* \sample \Kspace$
              \smallskip
    					\item $\advbit \gets \adv^{\RO}(\pk, \eK^*, \chalctxt)$
              \smallskip
    					\item $\pcoutput \bot$
    				\end{nicodemus}
    			\end{minipage}
    
    			\begin{minipage}[t]{5cm}
    				\underline{if $\adv$ calls $\RO(x)$}
            \smallskip
    				\begin{nicodemus}
              \item $c \assign \Enc_\pk(x)$ 
              \smallskip
              \item \pcif $c = \chalctxt :$ 
              \smallskip
              \item   \quad \green{\pcabort and \pcoutput $x$}
              \smallskip
              \item $\pcif \Db[x] = \bot :$ 
              \smallskip
              \item   \quad $\Db[x] \sample \bin^m$
              \smallskip
              \item $\pcreturn \Db[x]$
    				\end{nicodemus}	
    			\end{minipage}
    		}
    	\end{center}
    	\caption{\label{fig:cparom-red}
        Pseudocode for our reduction $\bdv$;
        highlighted in green is the point at which $\bdv$ will abort the
        simulation and win the one-wayness game. 
        If the simulation runs to completion without the correct value ever having
        been called to the random oracle, $\bdv$ returns $\bot$ to indicate 
        ``I give up''.
      }
    \end{figure}

    In summary: \emph{Either} $\adv$ never calls the
      random oracle on the encrypted value $x^*$ ($\bad = \false$), 
      in which case the view of $\adv$ is
      perfectly consistent with being in either of the real- or
      random-world games (advantage $= 0$), \emph{or} $\adv$
      calls $\RO$ on $x^*$ and therefore \emph{would} be able to
      distinguish ($\bad = \true$), in which case $\bdv$ 
      aborts the simulation and outputs $x^*$, 
      winning the one-wayness game.
    \qed
  \end{proof}

  %% IN CASE OF MULTI-CHAL VERSION
    %If this were all, then we would get a completely tight bound. 
    %  There is still one avenue of attack available to $\adv$, however,
    %  namely to guess the message underlying a \emph{future} challenge,
    %  then call the random oracle on that value before $\bdv$ has a
    %  chance of recognizing it as such. When that challenge later appears,
    %  $\adv$ will notice that the key $K$ it is issued in the challenge does
    %  not match the key that it had already computed.
    %
    %In Exercise~\ref{ex:remove-loss-term}, you are tasked with coming up with
    %  a better reduction that also handles this case. 
    %  %One solution would be to let our reduction check the query database for
    %  %collisions every time it receives a challenge, which would again make for
    %  %a tight reduction. 
    %  %in order to demonstrate techniques that will be useful later on when we move to the
    %  %quantum case, 
    %  Here, we opt for a ``dumber'' route, in order to demonstrate one of our
    %  main proof techniques,
    %  namely to hop to a game that simply aborts if this ``bad event'' ever
    %  occurs, and then argue that the probability this actually makes the game abort
    %  must be small. Then we can return to our main strategy, which will now 
    %  work without issue.

\subsection{When You Have to Guess}
  What would happen if we started from a probabilistic rather than a
    deterministic PKE? Well, assuming the probability of getting the same
    ciphertext twice is small (as is usually the case), then the re-encryption
    check that the reduction $\bdv$ relied on to find the correct pre-image,
    would in all likelihood fail to identify $x^*$ even if it were input.
    In other words, the reduction no longer has the ability to \emph{recognize} $x^*$, and
    so it must resort to \emph{guessing}.

  The argument goes: If $\adv$ is able to distinguish, then it must have
    queried the random oracle on input $x^*$ at \emph{some} point. 
    Therefore, with $\bdv$ keeping track of the oracle database $\Db$
    as part of its simulation, $x^*$ must be \emph{in} $\Db$, even if $\bdv$
    is unable to tell which one it is. So it does the simplest thing it can
    think of: After $\adv$ halts and outputs a guess $\advbit$, $\bdv$ chooses
    one of the $x$-values from $\Db$ uniformly at random, and returns that one
    as its guess. Then, if $\adv$ made $q$ queries to $\RO$, and provided
    $\bad = \true$ ($\adv$ called $\RO$ on $x^*$), $\bdv$ will win the
    one-wayness game with probability $\nicefrac{1}{q}$. After a little
    manipulation of probabilities, we end up with the following bound.

  \begin{theorem}[\PKE is \owcpa $\To$ \KEM is \indcpa]
    \label{thm:cparom-guess}
    Let \PKE be a probabilistic PKE scheme, let \RO be a random oracle, 
    and let the $\KEM = \U[\PKE, \RO]$.
    Then there is a simple, fully black-box reduction $\bdv$, such that for
    every adversary $\adv$,
    %that makes at most $\qe$ calls to the challenge oracle, 
    \[
      \advantage{\indcpa}{\KEM}[\left(\adv^{\RO}\right)] 
        \leq q \cdot \advantage{\owcpa}{\PKE}[\left(\bdv^{\adv}\right)]\, .
    \]
    The adversary $\bdv$ runs in the same time as $\adv$, plus the %(relatively small) 
    overhead of lazy-sampling a random oracle.
  \end{theorem}

  %\paragraph{A Loss of Tightness.}
  In \cref{thm:cparom}, the advantage of $\adv$ was \emph{tightly} bounded by
    the advantage of $\bdv$, in the sense that the advantage of $\bdv$ had a
    factor $1$ in front of it. Now there is a factor $q$. What is the
    consequence?

  Well, say we would like the advantage of $\adv$ to be upper-bounded by
    $2^{-128}$, a common target for security. If we start from a determinstic
    PKE scheme, we can rely on \cref{thm:cparom}, and choose parameters such
    that the best-known attack against the one-wayness of \dPKE is expected
    to succeed with probability $2^{-128}$ per time step -- or equivalently,
    that the best attacks require at least $2^{128}$ time steps in expectation
    to succeed. However, with $q$ counting the total number of hash function
    evaluations, $q$ can realistically get quite large: It is for example commonly
    estimated that the hash function SHA3 has been evaluated more than
    $2^{80}$ times in the life-time operation the Bitcoin block-chain. 

  Let $2^{-\lambda}$ denote the new one-wayness advantage that our parameters
    should target, after compensating for the loss in tightness. Inserting $q
    \leq 2^{80}$ gives us:
    \[
      \advantage{\indcpa}{\KEM}[\left(\adv^{\RO}\right)] 
        \leq q \cdot 2^{-\lambda} \leq 2^{80} \cdot 2^{-\lambda} = 2^{80 - \lambda} \leq 2^{-128}
        \implies \lambda \geq 208\, .
    \]
    Therefore, if we want the distinguishing advantage to be bounded by
      $2^{-128}$, we will need parameters such that the best attacks only
      achieve a one-way advantage of $2^{-208}$. 
      Perhaps not a dramatic increase --
      increasing the sizes of keys and ciphertexts accordingly should fairly
      easily be able to accomodate this -- but experience shows
      that even a modest loss in efficiency may make our scheme look a lot
      less competitive to implementers than than it just did, whether or not
      its competitors come bundled with tightness-respecting security proofs
      of their own!

    We will repeatedly tackle the question of tightness over the course of
      this tutorial, as one of the central challenges of the QROM is to pin
      down exactly what level of tightness is achievable, and when.

  \begin{exercise}\label{ex:cparom-lossy}
    Prove \cref{thm:cparom-guess}.\\
    \emph{
      Hint: Let $\game 1$ -- $\game 4$ be as before; you only need to change how
      the $\bad$ check is performed. Write down pseudocode for a
      reduction $\bdv$ that implements the guessing strategy described above, and
      lower-bound its one-wayness advantage in terms of $\prob{\bad}$ and $q$.
    }
  \end{exercise}

\subsection{One-way to Hiding}\label{sec:C-O2H}
  \paragraph{Find-O2H.}
    \Cref{thm:cparom} may be about a specific construction (the
      \U-transform), but the take-away message is quite general: If a reduction
      has control over the random oracle, and is able to recognize the correct
      answer when it sees it, we can use the random oracle to tightly reduce
      security from relying on a stronger distinguishing-style assumption, to
      instead rely on a weaker one-wayness assumption. 
      %If the reduction further has the ability to
      %recognize the correct answer when it sees it, this can be done tightly;
      %otherwise, the reduction will come with a loss factor $q$.
      We can generalize this observation in the following ``one-way to hiding'' lemma.\footnote{
        ``Hiding'' is another name for indistinguishability-style
        security, most often seen in the context of commitment schemes, which
        are usually required to be both ``hiding'' and ``binding''.
      }
      %One for the case that the reduction can \emph{find} the value it
      %Is looking for, and one for the case that it has no choice but to \emph{guess}.

    \begin{restatable}[Find-O2H, classical]{lemma}{cOWHfindsingle}\label{lemma:co2h-find-single}
      Let $\RO$, $\GO: \bin^n \to \bin^m$ be random
      oracles such that for every $x \neq x^*$, $\RO(x) = \GO(x)$, whereas
      $\RO(x^*)$ and $\GO(x^*)$ are independently uniform. 
      Let $z$ be a value that depends arbitrarily on $\RO$, $\GO$, and $x^*$.
      Let $\indfunc^*$ be an indicator oracle that on input $x$
      returns $1$ if $x = x^*$, and otherwise returns $0$. 

      Then there is a simple, fully black-box extractor $\extractor$,
      such that for every distinguisher $\distinguisher$,
      \[
          \abs{\prob{1 \gets \distinguisher^{\RO}(z)}
              - \prob{1 \gets \distinguisher^{\GO}(z)}}
            \leq \probsample{x \gets \extractor^{\distinguisher, \indfunc^*}(z)}{x = x^*}\, .
      \]
    \end{restatable}

    It may not be clear at a glance how this lemma relates to the proof of
      \cref{thm:cparom}. After all, there there was only one random oracle in
      play, while \cref{lemma:co2h-find-single} has two. Also,
      \cref{thm:cparom} speaks of an adversary and a reduction, while
      \cref{lemma:co2h-find-single} has a distinguisher and an extractors. How
      do they connect?

    To see the connection, note how the difference between $\game 2$ and $\game 3$
      is contained to the oracle interface only. In fact, the random oracle
      in $\game 2$ and the random oracle in $\game 3$ behave identically
      relative to $\adv$'s input,
      \emph{except} on the input value $x^*$, where it returns an independently
      uniform value instead of the expected $\eK^*$. 
      If we let the random oracle in $\game 2$ be identified with $\RO$, and the
      random oracle in $\game 3$ be $\GO$, and further let $z = (\pk, \eK^*,
      \ctxt^*)$, then we can identify $\distinguisher$ with $\adv$, and things
      start to line up:
      \[
        \abs{\probsample{\game{2}(\adv)}{1 \gets \adv}
          - \probsample{\game{3}(\adv)}{1 \gets \adv}}
        = \abs{\prob{1 \gets \distinguisher^{\RO}(z)}
          - \prob{1 \gets \distinguisher^{\GO}(z)}}\, .
      \]

    As for the extractor, it is given in \cref{fig:co2h-extractor}. It should
      look familiar: aside from the use of the ``indicator oracle'' $\indfunc^*$, it is an
      almost exact adaptation of $\bdv$ from \cref{fig:cparom-red}.
      The indicator oracle is simply a formalization of the statement that ``the
      extractor is able to recognize the correct value''; we provide this
      ability in oracle form to make sure that our formalization 
      only gives the extractor the ability to recognize correct values, 
      such that it can be applied to any situation where this is the case. For
      example, in our case the $\indfunc^*$ oracle can be locally computed from
      the input $z$, via the same re-encryption trick employed by $\bdv$.
      %and not accidentally give it an enhanced ability to find $x^*$ directly.

      %%% Extractor
      \begin{figure}[t]
      	\begin{center}
      		\fbox{
      			\small
      			\nicoresetlinenr
      			\begin{minipage}[t]{3.5cm}	
              \underline{\extHeading $\extractor^{\distinguisher, \indfunc^*}(z)$}
              \smallskip
      				\begin{nicodemus}
      					\item $\advbit \gets \distinguisher^{\RO}(z)$
                \smallskip
      					\item $\pcoutput \bot$
      				\end{nicodemus}
      			\end{minipage}
      
      			\begin{minipage}[t]{5cm}
      				\underline{if $\distinguisher$ calls $\RO(x)$}
              \smallskip
      				\begin{nicodemus}
                \item \pcif $\indfunc^*(x) = 1 :$ 
                \smallskip
                \item   \quad \pcabort and \pcoutput $x$
                \smallskip
                \item $\pcif \Db[x] = \bot :$ 
                \smallskip
                \item   \quad $\Db[x] \sample \bin^m$
                \smallskip
                \item $\pcreturn \Db[x]$
      				\end{nicodemus}	
      			\end{minipage}
      		}
      	\end{center}
      	\caption{\label{fig:co2h-extractor}
          The extractor $\extractor$ of \cref{lemma:co2h-find-single}.
          Here, the extractor $\extractor$ simulates $\RO$ for the distinguisher
          $\distinguisher$; the extractor might of course just as well provide a
          simulation of $\GO$.
        }
      \end{figure}

    Taken together, we can now provide simpler proof of the crucial hop
      from $\game 2$ to $\game 3$ in \cref{thm:cparom}, by identifying $\adv$
      with $\distinguisher$, applying \cref{lemma:co2h-find-single} to get the
      bound
      \[
        \abs{\probsample{\game{2}(\adv)}{1 \gets \adv}
          - \probsample{\game{3}(\adv)}{1 \gets \adv}}
        \leq \probsample{x \gets \extractor^{\distinguisher, \indfunc^*}(z)}{x = x^*}\, ,
      \]
      then let our reduction $\bdv$ instead simulate the environment for
      $\extractor$ by sampling $\eK^*$, setting $z = (\pk, \eK^*, \ctxt^*)$, 
      use the re-encryption trick to respond to $\indfunc^*$ queries, and output
      whatever value $\extractor$ outputs once it halts. Thus,
      \[
        \probsample{x \gets \extractor^{\distinguisher, \indfunc^*}(z)}{x = x^*}
        = \probsample{x \gets \bdv^{\adv, \extractor}(\pk, \ctxt^*)}{x = x^*}
        = \advantage{\owcpa}{\KEM}[(\bdv)]\, ,
      \]
      and the theorem is proven.

  \paragraph{Guess-O2H.}
    We can likewise define a ``Guess-O2H'' lemma, in correspondence with the guessing
      strategy of \cref{thm:cparom-guess}. Here the extractor is not given
      access to an indicator oracle, making the lemma applicable to situations
      where the reduction does not have the ability to recognize the correct
      answer. Hence, Guess-O2H has a higher level of applicability than Find-O2H, at the
      cost of sacrificing tightness.

    \begin{restatable}[Guess-O2H, classical]{lemma}{cOWHguesssingle}\label{lemma:co2h-guess-single}
      Let $\RO$, $\GO: \bin^n \to \bin^m$ be random
      oracles such that for every $x \neq x^*$, $\RO(x) = \GO(x)$, whereas
      $\RO(x^*)$ and $\GO(x^*)$ are independently uniform. 
      Let $z$ be a value that depends arbitrarily on $\RO$, $\GO$, and $x^*$.

      Then there is a simple, fully black-box extractor $\extractor$,
      such that for every distinguisher $\distinguisher$ that makes at most $q$
      calls to its random oracle,
      \[
          \abs{\prob{1 \gets \distinguisher^{\RO}(z)}
              - \prob{1 \gets \distinguisher^{\GO}(z)}}
            \leq q \cdot \probsample{x \gets \extractor^{\distinguisher}(z)}{x = x^*}\, .
      \]
    \end{restatable}

    \begin{exercise}\label{ex:cparom-lossy}
      Give a simpler proof of \cref{thm:cparom-guess} by
      referencing \cref{lemma:co2h-guess-single}. \\
      \emph{
        Hint: Let $\extractor$ do the heavy lifting -- no need to even specify
        how $\extractor$ works; it is enough that the lemma states that it
        exists. Then, all that is left for $\bdv$ to do is to provide $\extractor$ with its
        expected environment.
      }
    \end{exercise}

  %\subsubsection{Multi-point O2H.}
  %  The above lemmas (\ldots)
  %  \hhnote{cont}
  %\begin{lemma}[Find-O2H, classical]\label{thm:co2h-find}
  %  let $\RO: \bin^n \to \bin^m$ be a random
  %  oracle, and let $z$ be a bitstring that depends arbitrarily on $\RO$.
  %  Let $\solset \subseteq \bin^n$ be a subset of inputs, and let $\RO
  %  \backslash \solset$ be the punctured random oracle that is equal to $\RO$
  %  on all inputs $x \not\in \solset$, but which outputs $\bot$ on all points
  %  in $\solset$. Let $\find$ be the event that the random oracle ever outputs
  %  $\bot$. Then for every distinguisher $\distinguisher$,
  %  \[
  %      \abs{\prob{1 \sample \distinguisher^{\ro}(z)}
  %          - \prob{1 \sample \distinguisher^{\ro\backslash \solset}(z)}}
  %        = \probsample{\distinguisher^{\ro \backslash \solset}(z)}{\find}\, .
  %  \]
  %\end{lemma}


  %\hhnote{
  %  forward reference to QROM proof
  %}

  %\hhnote{todo: make exercises play well with autoref}
  %\begin{exercise}\label{ex:remove-loss-term}
  %  Write an alternative reduction that gives the same bound as in
  %  \cref{thm:cparom}, except without the extra term $\nicefrac{\qe}{2^m}$
  %  in the upper bound.
  %\end{exercise}

\subsection{The Post-Quantum Setting -- What Breaks?}
  Having seen the proofs of \cref{thm:cparom} and \cref{thm:cparom-guess}, 
    %\cref{lemma:co2h-find-single} and \cref{lemma:co2h-guess-single} 
    the above ``O2H lemmas'' may seem
    almost too trivial to bother writing down. After all, the extractors in the
    lemmas do little more than lazy-sample a random oracle, then either return
    the correct value when input (if the extractor is able to
    recognize it), or output one of the queries at random (if not).
    With the analysis being a near-trivial manipulation of
    conditional probabilities, it is not clear that the modularity gained 
    outweighs the cost of switching mental models from the behaviour of adversaries
    between games, to distinguishers trying to tell oracles apart.

  However, the same can certainly not be said for \emph{quantum} random
    oracles, since now, the input to the oracle can be a quantum superposition
    of any number of values at once. How does one recognize the correct input value
    in such a case? Is the ``bad'' event we used in our analysis
    even well-defined? Can you even \emph{guess}, given that the no-cloning
    theorem bars you from copying down the input states?

  This is where the very same O2H formalism that we introduced above earns its
    keep. In fact, the O2H lemmas turns out to be one the most useful toolkits
    we have, and they show up everywhere in the QROM literature. 
    This is because they successfully shove the tricky quantum analysis
    ``under the rug'', allowing non-quantum cryptographers to design their post-quantum
    game hops without having to worry that they are making a quantum fool of
    themselves. (Or at least, worry \emph{less}.)

  We will meet our first quantum O2H
    lemmas in \cref{sec:o2h}, where we will apply them to lift
    \cref{thm:cparom} and \cref{thm:cparom-guess} to the post-quantum
    setting. But before then, it is about time we brushed up on our quantum. 

  %\todoenum[rough structure]{
  %  \item A discussion about why the previous proofs do not yield
  %    post-quantum security 
  %  \item Where the proofs break down, and why they do not admit direct
  %    lifting (not history-free)
  %}
